\chapter{The geometry and field services}
\label{ch:services}

\chaptergoals{4 of 11 in Part~II.}
{\cref{ch:geometry} (\file{ship\_geometry.db}); the field maps from \cref{ch:setup}.}
{explain how other packages reach the geometry and the magnetic field, and inspect, evaluate and plot a field map.}

The geometry file and the field maps are just data. Two small libraries turn
them into something a simulation or reconstruction can use, without tying
either to a particular framework: neither library includes a Phlex header, so
these libraries can serve simulation, reconstruction and standalone programs.
The current \pkg{sea-cucumber} implementation opens GeoModelIO directly; its
optional geometry-service binding is still a stub.

\section{The geometry service}
\label{sec:geosvc}

\repo{geometry\_service} wraps the GeoModel description behind one abstract
interface, modelled on DD4hep's \texttt{IGeoSvc} (LHCb, Key4hep) and ALICE
O\textsuperscript{2}'s \texttt{GeometryManager}:

\begin{cpp}
namespace SHiP {
class IGeometryService {
 public:
  virtual ~IGeometryService() = default;
  virtual const GeoVPhysVol* geoModelWorld() const = 0;           // GeoModel tree
  virtual G4LogicalVolume& geant4WorldLogical() = 0;              // Geant4 world
  virtual G4LogicalVolume* getLogicalVolume(std::string_view name) const = 0;
};
}
\end{cpp}

\texttt{SHiP::GeometryService} implements it, with two ways to obtain the
geometry (before version~0.5 the class was \texttt{ship::SHiPGeometryService}):

\begin{cpp}
#include <SHiP/GeometryService/GeometryService.h>

// from a file written by build_geometry
auto svc = SHiP::GeometryService::fromFile("ship_geometry.db");
// or straight from the C++ description, without a file
auto svc2 = SHiP::GeometryService::fromSource();

const GeoVPhysVol* world = svc->geoModelWorld();
\end{cpp}

Two design decisions are worth knowing because they explain error messages
you may meet:

\begin{description}[style=nextline]
  \item[The GeoModel tree is built eagerly.] Both factory functions construct
    the full tree immediately, so a broken geometry fails at start-up, not in
    the middle of a run.
  \item[The Geant4 conversion is lazy.] \texttt{geant4WorldLogical()} converts
    GeoModel to Geant4 volumes on its first call only (guarded by
    \texttt{std::call\_once}), because the conversion needs a
    \texttt{G4RunManager} on the calling thread. Call it from
    \texttt{G4VUserDetectorConstruction::Construct()}. It returns the world
    by reference and throws if the conversion fails.
    \texttt{getLogicalVolume()} only works after that.
\end{description}

\cmd{ship\_g4view}, the viewer of \cref{sec:g4view}, is the smallest program
that does all this: \texttt{fromFile()}, then \texttt{geant4WorldLogical()}
inside \texttt{Construct()}, then one \texttt{G4PVPlacement} for the world.

The library also provides \texttt{G4RayScanner}, which walks a straight line
through the Geant4 geometry and returns the materials crossed with their
lengths (useful for material-budget studies), and \texttt{GeometryThread},
which funnels all Geant4 geometry creation through one thread, as Geant4's
multi-threading model requires.

\begin{exercise}{How aegir finds the geometry}
Open \file{aegir/src/geometry\_geomodel\_provider.cpp}. Without running
anything, answer:
\begin{enumerate}
  \item Which two configuration keys does the provider read?
  \item The workflow gives \texttt{db\_file: 'ship\_geometry.db'}. Where is
        the file looked for, and in which order?
  \item At which Phlex layer is the geometry published, and why does that
        make sense?
\end{enumerate}
\end{exercise}

\begin{solution}
\begin{enumerate}
  \item \texttt{db\_file} and \texttt{sensitive\_volumes}.
  \item First the path as given (so relative to the directory \cmd{phlex} runs
        in); if it does not exist, its bare file name in
        \file{\$SHIPGEOMETRY\_ROOT/share/geometry/}. Otherwise the provider
        throws ``Cannot locate geometry DB''.
  \item At the \texttt{job} layer: the detector does not change from event to
        event, so it is built once and shared by all events and threads.
\end{enumerate}
\end{solution}

\section{The field service}
\label{sec:fieldservice}

\repo{field\_service} stores magnetic field maps in a versioned file format and
evaluates them with \href{https://github.com/acts-project/covfie}{covfie}
(trilinear interpolation on a regular grid). It has four parts, so that each
consumer pulls in only what it needs:

\begin{center}
\begin{tabularx}{\linewidth}{l l L}
\toprule
Component & Needs & Purpose \\
\midrule
core       & --        & build a field evaluator from an in-memory map \\
MapIO      & ROOT      & read and write map files \\
G4Adapter  & Geant4    & a \texttt{G4MagneticField} for simulation \\
ActsAdapter& ACTS      & an \texttt{Acts::MagneticFieldProvider} for reconstruction \\
\bottomrule
\end{tabularx}
\end{center}

\subsection{The field-map format}

A field-map file is a ROOT file holding one RNTuple, \texttt{field\_maps},
with one entry per map, so a single file can carry several magnets. Each map
records its regular grid (first and last node and number of nodes per axis,
in mm), the three field components at every node (in tesla), the node
ordering, its symmetry, and where the values came from. Reading it needs
plain ROOT or uproot, not SHiP libraries. The full specification is
\file{docs/field\_map\_format.md} in the repository.

\paragraph{Symmetry.} A map can be stored for one quadrant only and mirrored
across the $x=0$ or $y=0$ planes, with a sign (\emph{parity}) per field
component and plane. FairShip's quadrant-symmetric muon-shield maps are stored
this way.

\paragraph{Frames.} A map is stored in the magnet's own frame. Whoever uses it
places it in the global frame with a translation. In \pkg{aegir}, each magnet is
also attached to the logical volumes whose names contain a given pattern;
Geant4 evaluates the field only inside those volumes, so field-free drift
regions cost nothing:

\begin{jsonnetcode}
field: lib.field_map([
  { name: 'MuonShield', volume_pattern: '/SHiP/muon_shield',
    file: 'muon_shield_TRY_2026.root', map: 'muon_shield_TRY_2026',
    translation: [0, 0, 2040] },     // shield entrance [mm]
  { name: 'Spectrometer', volume_pattern: '/SHiP/magnet',
    file: 'spectrometer_V21_2455.root', map: 'spectrometer_V21_2455',
    translation: [0, 0, 89570] },    // magnet centre [mm]
]),
\end{jsonnetcode}

\subsection{The command-line tools}

\begin{center}
\small
\begin{tabularx}{\linewidth}{l L}
\toprule
\cmd{fieldmap\_dump f.root}            & print the metadata of every map in the file \\
\cmd{fieldmap\_dump f.root map}        & read \texttt{x y z} lines (mm) from standard input,
                                         print \texttt{x y z Bx By Bz} (mm, T) \\
\cmd{plot\_field\_map f.root map out.root} & plot $B_y$ along the beam axis and in the
                                         $xz$, $yz$ and central $xy$ planes, as ROOT histograms
                                         and one PDF per plot \\
\cmd{generate\_constant\_fieldmap} & write a uniform field over a box, for tests \\
\cmd{fairship\_to\_fieldmap}       & convert a FairShip map; \cmd{--symmetry quadrant\_dipole}
                                     for the quadrant-symmetric ones \\
\bottomrule
\end{tabularx}
\end{center}

The tools are in the \texttt{field-service-tools} conda package, which the
\repo{field\_maps} workspace of \cref{ch:setup} uses to convert the maps; its
environment is where you run them.

In this tutorial the muon-shield map matters most: its field is what sweeps
the muons from the target away from the decay volume in \cref{ch:aegir}.

\begin{exercise}{Inspect the field maps}
In the \repo{field\_maps} environment, in \file{work/}, where the maps are:
\begin{shell}
cd "$SHIP_TUTORIAL/field_maps"
pixi cvmfs
cd "$SHIP_TUTORIAL/work"
fieldmap_dump spectrometer_V21_2455.root
fieldmap_dump muon_shield_TRY_2026.root
\end{shell}
For each map, note the grid extent and spacing, and whether it is mirrored.
\end{exercise}

\begin{exercise}{Evaluate and plot the spectrometer field}
The spectrometer map is stored in the magnet's own frame, so its centre is
at $(0,0,0)$. Evaluate the field there and one metre either side along the
beam:
\begin{shell}
printf '0 0 0\n0 0 -1000\n0 0 1000\n' | \
  fieldmap_dump spectrometer_V21_2455.root spectrometer_V21_2455
plot_field_map spectrometer_V21_2455.root spectrometer_V21_2455 \
  spectro_plots.root
\end{shell}
Which component dominates? Open the $B_y$-versus-$z$ PDF and estimate the
field integral $\int B_y\,\mathrm{d}z$. The transverse momentum kick is
$\Delta p_T\,[\mathrm{GeV}] \approx 0.3 \int B\,\mathrm{d}z\,[\mathrm{T\,m}]$:
what does that mean for the bending of a 5\,GeV muon crossing the spectrometer?
\end{exercise}

\begin{pitfall}
Outside its grid the evaluator clamps to the boundary value. A flat plateau
beyond the edges of a plot is an artefact of the plotting, not a physical
fringe field. The same clamping applies in the simulation, which evaluates the
field only inside the volumes matching \texttt{volume\_pattern}: check that the
map's grid covers those volumes, or Geant4 sees the boundary value continued
up to the volume edge.
\end{pitfall}

\begin{note}
For a demonstration, a uniform spectrometer field can replace the real map. For example, 1\,T vertical over the magnet volume:
\begin{shell}
generate_constant_fieldmap spectrometer_uniform.root spectrometer_uniform \
  0 1 0  -2500 2500  -3500 3500  -2500 2500
\end{shell}
(arguments: file, map name, $B_x\,B_y\,B_z$ in T, then
$x_\mathrm{min}\,x_\mathrm{max}\,y_\mathrm{min}\,y_\mathrm{max}\,z_\mathrm{min}\,z_\mathrm{max}$
in mm). To use it, open a copy of \file{muon\_gun.jsonnet} in your editor
and set \texttt{file} and \texttt{map} of the \texttt{Spectrometer} entry to
\texttt{spectrometer\_uniform.root} and \texttt{spectrometer\_uniform}.
The uniform-field sample is not suitable for an absolute performance prediction.
\end{note}

\begin{organisernote}
The maps are those of \repo{field\_maps}: FairShip's \file{TRY\_2026.root}
and \file{2026\_07\_02\_MainSpectrometerField\_V21\_2455.root} at the FairShip
commit pinned in its \file{placements.json}, converted with
\cmd{fairship\_to\_fieldmap} (\cmd{--symmetry quadrant\_dipole} for the
muon shield). Its \file{field\_maps.libsonnet} gives the same magnets and
translations as the workflows of this booklet. The muon-shield grid is
$111\times110\times651$ nodes over $x,y\in[0,2.2]$\,m,
$z\in[-0.5,32]$\,m; the spectrometer grid $101\times101\times151$ nodes over
$x,y\in[-5,5]$\,m, $z\in[-7.5,7.5]$\,m. Fill in the expected numbers for
the exercise: $B_y$ at the centre and the field integral.
\end{organisernote}
